\subsection*{الف}

تبدیل فوریهٔ $y^*[-n]$ برابر است با $Y^*(e^{j\omega})$.
همچنین می‌دانیم ضرب در حوزهٔ فرکانس متناظر با کانولوشن در حوزهٔ زمان است.

پس اگر:
\[
G(e^{j\omega}) = X(e^{j\omega}) Y^*(e^{j\omega})
\]
آنگاه در حوزهٔ زمان خواهیم داشت:
\[
g[n] = x[n] * y^*[-n]
\]

بنابراین این دو یک زوج تبدیل فوریه تشکیل می‌دهند.

\subsection*{ب}

از قسمت (الف) داریم:
\[
g[n] = x[n] * y^*[-n]
\quad \Longleftrightarrow \quad
G(e^{j\omega}) = X(e^{j\omega})Y^*(e^{j\omega})
\]

اکنون از رابطهٔ تبدیل فوریهٔ معکوس استفاده می‌کنیم:
\[
g[n] =
\frac{1}{2\pi}
\int_{-\pi}^{\pi}
X(e^{j\omega})Y^*(e^{j\omega}) e^{j\omega n} d\omega
\]

از طرف دیگر، با باز کردن کانولوشن:
\[
g[n] = \sum_{k=-\infty}^{\infty} x[k]\, y^*[k-n]
\]

پس:
\[
\frac{1}{2\pi}
\int_{-\pi}^{\pi}
X(e^{j\omega})Y^*(e^{j\omega}) e^{j\omega n} d\omega
=
\sum_{k=-\infty}^{\infty} x[k]\, y^*[k-n]
\]

حال اگر $n=0$ قرار دهیم:
\[
\frac{1}{2\pi}
\int_{-\pi}^{\pi}
X(e^{j\omega})Y^*(e^{j\omega}) d\omega
=
\sum_{k=-\infty}^{\infty} x[k]\, y^*[k]
\]
